\section{Neural LoFi from Gradient Descent}

\label{2605:app:GD-neural}
In this section, our goal is to prove Proposition \ref{2605:thm:layerwise-GD}, whose formal version is stated below as Theorem \ref{2605:thm:neural-lofi}, i.e. show that under layerwise training with vanishing initialization, the early-time dynamics of each layer produces spikes along the linear direction $\hat{\bu}_\ell$ and aligns with the top eigenvectors of $\widehat{C}^{(\ell)}$, for a standard fully-connected architecture without skip connections.

Consider the standard $L$-layer network
\begin{equation}
\hat f(\bx)=\langle a_L, \bz_L(\bx)\rangle,
\qquad
\bz_0(\bx)=\bx,
\qquad
z_\ell(x)=\sigma\!\bigl(W_\ell \bz_{\ell-1}(\bx)\bigr),\quad \ell=1,\dots,L,
\label{2605:eq:sec2-arch}
\end{equation}
with square loss $\mathcal{L}=\tfrac{1}{2n}\sum_\mu(y_\mu-\hat f(\bx_\mu))^2$. The final readout $a_L\in\mathbb{R}^{p_L}$ is fixed throughout training; the weight matrices $W_\ell$ are updated layer-wise.




We write $w_{\ell,i}$ for the $i$-th row of $W_\ell$ and set
\begin{equation}\label{2605:eq:init-scales}
\alpha_m := \max_i \|\bw_{m,i}(0)\|, \qquad m=1,\dots,L,
\end{equation}
with  $\alpha_{L+1}:=\|a_L\|$. 

\begin{assumption}[Layerwise scale separation]\label{2605:ass:scale-sep}
The dataset, architecture, and activation are fixed. The scales, including the
final readout scale \(\alpha_{L+1}=\|a_L\|\), form a strict hierarchy: for
every $\ell$ and every $m\in\{\ell+1,\dots,L+1\}$,
\begin{equation}\label{2605:eq:scale-hierarchy}
\alpha_m=o(\alpha_\ell).
\end{equation}
\end{assumption}

Thus later hidden layers, and the final readout, vanish on a strictly smaller
asymptotic scale than the layer currently being trained.  The hierarchy is used
below in a scale-counting sense: every replacement of a later-layer quantity by
its leading initialization-scale term carries at least one additional factor
from some \(\alpha_m\) with \(m>\ell\).

\begin{assumption}\label{2605:ass:nonlin}
$\sigma:\mathbb{R}\to\mathbb{R}$ is three times continuously differentiable,
$\sigma(0)=0$, and $\|\sigma'\|_\infty,\|\sigma''\|_\infty,\|\sigma'''\|_\infty<\infty$.
\end{assumption}

Define the \emph{linear spike direction} and the \emph{label-weighted covariance} at layer $\ell$:
\begin{equation}\label{2605:eq:sec2-linear-component}
\hat{\bu}_\ell
\coloneqq \frac{1}{n}\sum_{\mu=1}^n y_\mu\,\bz_{\ell-1}(\bx_\mu),
\qquad
\widehat{C}^{(\ell)}
\coloneqq \frac{1}{n}\sum_{\mu=1}^n
y_\mu\,\bz_{\ell-1}(\bx_\mu)\bz_{\ell-1}(\bx_\mu)^\top.
\end{equation}

\begin{theorem}[Layerwise GD approximation]\label{2605:thm:neural-lofi}
Under Assumptions~\ref{2605:ass:scale-sep}--\ref{2605:ass:nonlin}, there exist constants
$\tau,C>0$ such that the layerwise gradient descent run admits
horizons
\begin{equation}\label{2605:eq:layerwise-horizons}
T_\ell:=\left\lfloor \frac{\tau\alpha_\ell}{\eta}\right\rfloor,
\qquad \ell=1,\dots,L,
\end{equation}
with the following property: for every layer $\ell$, every neuron $i$, and every
$0\le t\le T_\ell$,
\[
\|w_{\ell,i}(t)\|\le C\alpha_\ell.
\]
Thus the layer receives an $O(\alpha_\ell)$ update while staying on its
initialization scale.  Define the leading effective readout
\begin{equation}\label{2605:eq:effective-readout-fc}
\bar{a}_{\ell,i}
:= c_0^{L-\ell}\Bigl[\Bigl(\prod_{m=\ell+1}^{L}W_m(0)\Bigr)^\top a_L\Bigr]_i
\end{equation}
where \(c_0=\sigma'(0)\). Suppose that the effective readout for neuron $i$ is non-degenerate i.e.:
\begin{equation}\label{2605:eq:readout-nondegenerate}
    |\bar a_{\ell,i}|\ge c_{\rm rd} \alpha_{L+1}\prod_{m=\ell+1}^{L}\alpha_m
\end{equation}
for a fixed \(c_{\rm rd}>0\). For such neurons and \(0\le t<T_\ell\),
\begin{equation}\label{2605:eq:sec2-main-dynamics}
w_{\ell,i}(t+1) - w_{\ell,i}(t)
= \eta\,c_0\,\bar{a}_{\ell,i}\,\hat{\bu}_\ell
+ \eta\,\bar{a}_{\ell,i}\,c_1\,\widehat{C}^{(\ell)}\,w_{\ell,i}(t)
+ R_{\ell,i}(t),
\end{equation}
where \(c_1=\sigma''(0)\). For every \(0\le t\le T_\ell\) the accumulated residual obeys
\begin{equation}\label{2605:eq:remainder-bound}
    \left\|
        \sum_{s=0}^{t-1}R_{\ell,i}(s)
    \right\|
    \le
    C\alpha_\ell^3 .
\end{equation}
Equivalently, the spike contributes \(O(\alpha_\ell)\) over the horizon, the
covariance-amplification term contributes \(O(\alpha_\ell^2)\), and the
discarded part is \(O(\alpha_\ell^3)\). Thus the displayed approximation keeps
the first two nontrivial orders in the layer-\(\ell\) initialization scale.
\end{theorem}

\begin{remark}
The three components in Equation~\eqref{2605:eq:sec2-main-dynamics} have the following roles:
\begin{itemize}[noitemsep]
\item \textbf{Linear spike.} $\eta c_0 \bar{a}_{\ell,i}\hat{\bu}_\ell$ is independent of the current weight; over the horizon $T_\ell$ it moves the neuron by $O(\alpha_\ell)$ toward the empirical label-feature correlation.
\item \textbf{Covariance amplification.} $\eta \bar{a}_{\ell,i} c_1\widehat{C}^{(\ell)}w_{\ell,i}(t)$ is linear in the current weight and amplifies components aligned with the leading eigenvectors of $\widehat{C}^{(\ell)}$.
\item \textbf{Remainder.} The per-step remainder is second order in the current
layer scale, hence it accumulates to \(O(\alpha_\ell^3)\) over
\(T_\ell\asymp \alpha_\ell/\eta\) steps. Terms coming from later layers are
smaller by the scale hierarchy in Assumption~\ref{2605:ass:scale-sep}.
\end{itemize}
The proof below shows that \(\bar a_{\ell,i}\) is the leading, sample-independent
backpropagation coefficient along the path from layer \(\ell\) to the readout.
The non-degeneracy condition only excludes neurons whose leading
backpropagation coefficient is accidentally cancelled; such neurons are inactive
at the displayed order.
\end{remark}

\subsection{Proof of Theorem~\ref{2605:thm:neural-lofi}}

\begin{proof}
All constants below may depend on the fixed data and architecture bounds,
$L$ and the derivative bounds of $\sigma$.
We write \(K\) for a finite constant of this kind, changing from line to line,
and do not track separate data, width, or label constants.

For the standard architecture the gradient of the loss w.r.t.\
\(w_{\ell,i}\) involves a single backpropagation path through all layers
\(\ell+1,\dots,L\). Define the layer-to-layer Jacobian
\begin{equation}\label{2605:eq:jacobian}
\bm{J}_{L\ell}(x)
:= \frac{\partial z_L(x)}{\partial z_\ell(x)}
= \prod_{m=\ell+1}^{L}\mathrm{diag}\!\bigl(\sigma'(\bm{h}_m(x))\bigr)W_m,
\quad \bm{h}_m(x)=W_m z_{m-1}(x),
\end{equation}
and the sample-dependent effective readout
\[
    \bar a_{\ell,i}(x):=[\bm J_{L\ell}(x)^\top a_L]_i .
\]
By the chain rule,
\begin{equation}\label{2605:eq:full-grad}
\nabla_{w_{\ell,i}}\mathcal{L}
= -\frac{1}{n}\sum_{\mu=1}^n
  r_\mu\,\bar{a}_{\ell,i}(\bx_\mu)\,
  \sigma'\!\bigl(u_{\ell,i}(\bx_\mu)\bigr)\,\bz_{\ell-1}(\bx_\mu),
\quad r_\mu := y_\mu - \hat f(\bx_\mu),
\end{equation}
where $u_{\ell,i}(x)=\langle w_{\ell,i},\bz_{\ell-1}(x)\rangle$.

The key observation is that, under Assumption~\ref{2605:ass:scale-sep}, the activation
derivatives along this path are well approximated by their value at zero, so
\(\bar a_{\ell,i}(x)\) is approximated by the initialization-dependent
coefficient in~\eqref{2605:eq:effective-readout-fc}. On the layerwise horizons
considered below the representations are uniformly bounded and
\(\|W_m\|_{\mathrm{op}}\lesssim \alpha_m\). Hence
\(\|\bm h_m(x)\|\lesssim\alpha_m\), and Taylor expanding \(\sigma'\) around
zero gives
\begin{equation}\label{2605:eq:diag-expand}
\mathrm{diag}(\sigma'(\bm{h}_m(x))) = c_0 I + c_1\,\mathrm{diag}(\bm{h}_m(x)) + O(\alpha_m^2),
\quad c_0=\sigma'(0),\;c_1=\sigma''(0).
\end{equation}
Let
\[
    \bm J_{L\ell}^{(0)}:=c_0^{L-\ell}\prod_{m=\ell+1}^L W_m(0)
\]
with the convention that the empty product is the identity. Expanding the
product for \(\bm J_{L\ell}(x)\) around this leading term gives the explicit
scale bound
\[
    \|\bm{J}_{L\ell}(x)-\bm J_{L\ell}^{(0)}\|_{\mathrm{op}}
    \le
    K
    \left(\prod_{m=\ell+1}^{L}\alpha_m\right)
    \left(\sum_{m=\ell+1}^{L}\alpha_m\right).
\]
Since each \(\alpha_m=o(\alpha_\ell)\) for \(m>\ell\), the right-hand side is
\(o(\alpha_\ell)\). Therefore
\begin{equation}\label{2605:eq:Jac-correction}
\|\bm{J}_{L\ell}(x) - \bm{J}_{L\ell}^{(0)}\|_{\mathrm{op}}
= o(\alpha_\ell),
\end{equation}
and when \(\ell=L\) the left-hand side is zero. Multiplying by the final
readout gives, with \(\gamma_\ell=\alpha_{L+1}\prod_{m=\ell+1}^L\alpha_m\),
\[
    |\bar a_{\ell,i}(x)-\bar a_{\ell,i}|
    \le
    K\gamma_\ell
    \left(\sum_{m=\ell+1}^{L}\alpha_m\right),
\]
while \eqref{2605:eq:readout-nondegenerate} gives the relative estimate
\begin{equation}\label{2605:eq:readout-const-approx}
    \frac{|\bar{a}_{\ell,i}(x)-\bar{a}_{\ell,i}|}
    {|\bar a_{\ell,i}|}
    \le
    C
    \sum_{m=\ell+1}^{L}\alpha_m
    =
    o(\alpha_\ell).
\end{equation}
In particular,
\(\bar a_{\ell,i}(x)=\bar a_{\ell,i}(1+\varepsilon_{\ell,i}(x))\) with
\(\sup_x|\varepsilon_{\ell,i}(x)|=o(\alpha_\ell)\). The error is zero when \(\ell=L\).

We next translate the above bound to the required bound on the dynamics of $W_\ell$. We start by stating uniform bounds required throughout.  Since the sample,
widths, and activation are fixed, and since we may restrict to scales
\(\max_m\alpha_m\le1\), the row-norm condition
\(\|w_{j,i}\|\le2\alpha_j\) implies recursively that
\[
    \max_{\mu,j}\|z_j(\bx_\mu)\|\le K .
\]
Hence, at the start of
layer \(\ell\),
\[
\|\hat{\bu}_\ell\|
\le \frac1n\sum_\mu |y_\mu|\,\|\bz_{\ell-1}(\bx_\mu)\|
\le K,
\]
and
\[
\|\widehat C^{(\ell)}\|_{\mathrm{op}}
\le \frac1n\sum_\mu |y_\mu|\,\|\bz_{\ell-1}(\bx_\mu)\|^2
\le K .
\]

The same bounds give \(\|\bm J_{L\ell}^{(0)}\|_{\mathrm{op}}\le K\) and hence
uniformly bounded effective readouts.  The residuals \(r_\mu\) are uniformly
bounded as well.  Therefore
\(\|\nabla_{w_{\ell,i}}\mathcal L\|\le K\).  Choose \(\tau\) small enough that
\(\tau K\le 1\).  Then, for every
$t$ with $\eta t\le \tau\alpha_\ell$,
\[
\|w_{\ell,i}(t)\|
\le \|w_{\ell,i}(0)\|+\eta tK
\le 2\alpha_\ell .
\]

On this horizon,
\(|u_{\ell,i}(\bx_\mu)|\le K\alpha_\ell\).  Taylor's theorem and
Assumption~\ref{2605:ass:nonlin} give
\begin{equation}\label{2605:eq:sec2-sigma-linearization}
\sigma'(u) = c_0 + c_1\,u + \omega_\sigma(u),\quad
|\omega_\sigma(u)|\leq \tfrac{1}{2}\|\sigma'''\|_\infty u^2 .
\end{equation}
Moreover, by~\eqref{2605:eq:readout-const-approx},
\(\bar a_{\ell,i}(\bx_\mu)=\bar a_{\ell,i}(1+\varepsilon_{\ell,i}(\bx_\mu))\)
with \(\sup_\mu|\varepsilon_{\ell,i}(\bx_\mu)|=o(\alpha_\ell)\).  Thus the
readout replacement error is smaller than the leading readout contribution
itself; quantitatively it is
\(|\bar a_{\ell,i}|\,o(\alpha_\ell)=o(\alpha_\ell^2)\), because
\(|\bar a_{\ell,i}|=O(\gamma_\ell)=o(\alpha_\ell)\).  By the hierarchy of scales, we further have
$|\hat f(\bx_\mu)|=o(\alpha_\ell^2)$, and
$r_\mu=y_\mu+o(\alpha_\ell^2)$.
Substituting these three estimates in~\eqref{2605:eq:full-grad} yields
\[
\frac1n\sum_{\mu=1}^n
r_\mu\,\bar a_{\ell,i}(\bx_\mu)\,
\sigma'\!\bigl(u_{\ell,i}(\bx_\mu)\bigr)\,\bz_{\ell-1}(\bx_\mu)
=
\bar a_{\ell,i}\bigl[c_0\hat{\bu}_\ell
+c_1\widehat C^{(\ell)}w_{\ell,i}(t)\bigr]
+\mathcal E_{\ell,i}(t),
\]
with the per-step error bound
\[
\|\mathcal E_{\ell,i}(t)\|
\le
K\Bigl(
    |\bar a_{\ell,i}|\,\sup_\mu|\varepsilon_{\ell,i}(\bx_\mu)|
    + |\bar a_{\ell,i}|\alpha_\ell^2
    + o(\alpha_\ell^2)
\Bigr)
\le K\alpha_\ell^2 .
\]
Finally,
$w_{\ell,i}(t+1)-w_{\ell,i}(t)=-\eta\nabla_{w_{\ell,i}}\mathcal L$, so setting
$R_{\ell,i}(t):=\eta\mathcal E_{\ell,i}(t)$ proves
\eqref{2605:eq:sec2-main-dynamics}.  Summing over \(0\le s<t\le T_\ell\) gives
\[
    \left\|
        \sum_{s=0}^{t-1}R_{\ell,i}(s)
    \right\|
    \le
    \eta t\,K\alpha_\ell^2
    \le
    K\alpha_\ell^3,
\]
which is \eqref{2605:eq:remainder-bound}.
\end{proof}
